% EN49 c02 | 原书 p.12–21（PDF 23–32） | 第 1 章 §1.4–1.5
% 骨架文件：标题行已设定，正文由译者填写。每页起首用 % ===== 原书 p.N ===== 注释标记。
\subsection{大复结构极限处的镜像对称}

% ===== 原书 p.12 =====
上文给出的奇异超曲面的深谷范畴的定义，显然依赖于磨光（smoothing）的选择。这不仅是合意的，而且对同调镜像对称而言至关重要。经典地，镜像对称是凯勒流形与大复结构极限（large complex structure limit，LCSL）复流形族之间的关系。同调镜像对称被预期为一个对合，因此有必要为该族的奇异纤维的深谷范畴建立一个概念。既然镜像的选择依赖于整个退化，奇异纤维的深谷范畴之定义应当涉及磨光的数据，也就不足为奇了。

\transnote{易错预警：原文说同调镜像对称“被预期为一个对合”（expected to be an involution），意思是镜像对称操作两次应回到原处，而不是说它作为函子是自同构。这一预期正是必须为奇异纤维本身定义深谷范畴的理由：若只在邻近纤维上定义，A 侧就无法与 B 侧对称地取镜像。}

另一方面，孤立超曲面奇点的芽具有光滑且连通的 versal 形变空间 [KM98, p. 144]，故一个简单的同伦论证表明：此时（芽的）深谷范畴与磨光的选择无关。因此它给出孤立超曲面奇点的一个（可能有趣的）辛不变量。一般情形下必须提供额外数据，深谷范畴才能被唯一确定：§5 给出了例子。来自镜像对称的预期表明，这一额外数据应采取 $f^{-1}(0)$ 上对数结构（log structure）的形式，因为在好的情形下它被预期决定 $f^{-1}(0)$ 的一个磨光。利用 Parker 的爆炸流形（exploded manifolds）上全纯曲线理论，或许可以为带对数结构的奇异超曲面内蕴地构造 wrapped Fukaya 范畴（某些对数概型具有爆炸流形的结构 [Par12]）。

\transnote{含糊点拨：对数结构（log structure）在这里只是一个来自镜像对称的\emph{预期}，原文并未给出其定义或构造。它大致上记录边界除子与极点的数据，从而补充“沿哪个方向磨光”的信息；把它当作“把磨光唯一确定下来所需的额外数据”来读即可。}

% ===== 原书 p.13 =====
换个角度，从 LG 模型 $(X\times\C,zf)$ 来看，临界轨迹 $f^{-1}(0)\times\C$ 带有 $(-1)$-平移辛结构（$(-1)$-shifted symplectic structure）：我们猜想这一额外结构也足以确定 wrapped Fukaya 范畴，而后者可以用 Joyce 的 $d$-临界轨迹（$d$-critical loci）理论 [Joy15] 等框架内蕴地构造。

\transnote{易错预警：“$(-1)$-平移辛结构”中的平移量 $(-1)$ 是移位辛几何里的移位次数，与“负一次”或“倒数”无关；不要读成辛形式的次数为 $-1$。$d$-临界轨迹的 $d$ 同样是理论的名字（Joyce 的 $d$-critical loci），不是微分符号。}

当辛侧是形如 $(X\times\C,zf)$ 的 LG 模型时，镜像对称的结论可借定理 \ref{thm:knorrer} 改写成关于奇异簇 $f^{-1}(0)$ 的镜像对称陈述。例如，Nadler 的两篇论文 [Nad19a; Nad17] 研究了 $(\C^{n+1},z_1\cdots z_{n+1})$ 的 A 模型的微局部版本。在此情形下，我们的周期性定理有另一种证明，我们将在 §5.1 中概述；此外还可以用 Abouzaid–Auroux 的工作 [AA21]，以我们的术语给出 Nadler 结果的一个简单证明。我们借用即将问世的 [AA21] 续篇中的想法，说明我们的定义如何使这一点推广到环簇中超曲面的补集的情形。

这是更一般的镜像对称陈述的一个特例，后者可在 Gross–Siebert 纲领的框架下理解（记号见 §5.3）。设 $B$ 是整仿射流形，带有多面体分解 $\cP$ 以及严格凸整分段线性多值函数 $\varphi$ 的选择，并令 $(\check B,\check\cP,\check\varphi)$ 为离散 Legendre 对偶数据。再令 $\cX_\varphi\to D$ 与 $\check\cX_{\check\varphi}\to D$ 分别是 Gross 与 Siebert 构造的 $X(B)=TB/T^{\Z}B$ 与 $X(\check B)=T\check B/T^{\Z}\check B$ 的环簇退化。大体积极限是补集 $X(B)\setminus s^{-1}(0)$，其中 $s$ 是联系于 $\varphi$ 的丰富线丛 $\cL_\varphi$ 的一个截面；而大复结构极限是某个环簇退化的中心纤维 $\cX^0_\varphi$。

\begin{theorem}[{大复结构极限与大体积极限的镜像对称}]\label{thm:lvl-lcsl}
当 $B=M_\R/\Gamma$ 是整仿射环面，带有多面体分解 $\cP$，并取定严格凸分段线性多值函数 $\varphi$（其斜率为整）时，大复结构极限 $\check\cX^0_{\check\varphi}$ 同调镜像于大体积极限 $X(B)\setminus s^{-1}(0)$，即存在范畴的拟等价
\[
\mathrm{D}^\pi\cF(\check\cX^0_{\check\varphi})\simeq\mathrm{D}^b\operatorname{Coh}(X(B)\setminus s^{-1}(0)).
\]
\end{theorem}

\transnote{易错预警：花体 $\cX_\varphi$、$\check\cX_{\check\varphi}$ 记环簇退化的\emph{全族}，罗马体 $X(B)$、$X(\check B)$ 记其一般纤维；中心纤维是 $\cX^0_\varphi$（上标 $0$，花体），不能与 $X(B)$ 混同。定理中的 $D$ 是退化族的底，不是除子。}

% ===== 原书 p.14 =====
该定理的一个推论是：带 $n$ 个结点的椭圆曲线的 Fukaya 范畴，与带 $n$ 个穿孔的椭圆曲线的导出相干层范畴之间，存在一个镜像对称陈述。据我们所知，这些同调镜像对称的实例此前并不为人所知。更一般地，我们有

\begin{theorem}[{大复结构极限处的镜像对称}]\label{thm:lcsl}
给定 $X(B)$ 与 $X(\check B)$ 之间的一个同调镜像对称等价，它把绕 $X(\check B)$ 的大复结构极限的（顺时针）单值函子 $\mu$ 与 $X(B)$ 上线丛 $\cL^{-1}$ 的张量积等同起来，且 $\cL$ 的第一 Chern 类 $c_1(\cL)$ 就是 $X(B)$ 上凯勒形式所在的类；则 $X(\check B)$ 的大复结构极限 $X_0(\check B)$ 同调镜像于 $X(B)$ 的某个大体积极限：
\[
\mathrm{D}^\pi\cF(X_0(\check B))\simeq\mathrm{D}^b\operatorname{Coh}(X(B)\setminus s^{-1}(0)),
\]
其中 $s$ 是 $\cL$ 的截面。
\end{theorem}

\transnote{易错预警：定理 \ref{thm:lcsl} 中的大复结构极限写作 $X_0(\check B)$（罗马体 $X$ 带下标 $0$），与定理 \ref{thm:lvl-lcsl} 的花体 $\check\cX^0_{\check\varphi}$ 指同一类对象，只是记号不同。原文如此，照译，不要自行统一。另：$X(B)$ 与 $X(\check B)$ 是互为镜像的一对环簇，$\check B$ 上的数据才对应 $B$ 侧，方向容易读反。}

应用定理 \ref{thm:lcsl} 所需形式的同调镜像对称等价，被预期可由 Abouzaid 的族 Floer 理论结合 Gross–Siebert 纲领的想法推出（见 [GS16]）。§5.3 给出了进一步解释：这些论证必然带有某种启发性质，忽略了局部系统、brane 结构等要点。

% ===== 原书 p.15 =====
等等。定理 \ref{thm:lcsl} 应被当作一个例证：定义 \ref{def:sing-fuk} 应当给出所预期的那些镜像对称等价。

\begin{remarkx}\label{rem:lvl-direction}
证明同调镜像对称的许多策略，与定理 \ref{thm:lcsl} 的方向相反：先证明大复结构极限（LCSL，在 B 侧）与大体积极限（LVL，在 A 侧）之间的等价，再用形变理论的论证推出邻近纤维的同调镜像对称。虽然理解 LVL 的深谷范畴比理解邻近纤维的更容易，但 LCSL 是奇异簇，其深谷范畴不能按常规方式定义，只能借助邻近纤维来定义。这为两种进路之间的差别提供了一些理由；当然，可以把两种进路结合起来，以获得对模空间上各点处同调镜像对称的理解。
\end{remarkx}

\transnote{背景直觉：注 \ref{rem:lvl-direction} 点出了本论文与既有文献的分工。传统路线是“先在极限处证等价，再形变到一般纤维”；本文的路线是把奇异纤维的深谷范畴本身定义好，从而让“极限处的等价”就是一个关于奇异簇的陈述。两条路线合起来覆盖整个模空间。}

\subsection{环簇的镜像对称}

本论文的许多结果都在环簇 Fano 簇的镜像对称背景下展开，这里先作回顾。设 $H\subseteq(\C^*)^n$ 是光滑代数超曲面（即极仿射超曲面），并取 $n$ 元 Laurent 多项式 $f$，使 $H=f^{-1}(0)$。我们将 $f$ 写成如下形式
\[
f(z)=\sum_{\alpha\in A}c_\alpha T^{\varphi(\alpha)}z^\alpha
\]
其中 $A\subseteq\Z^n$ 是 $f$ 中所有单项式的多重指标之集；$c_\alpha$ 暂时取为任意非零复数；$T$ 是一个小实数；$\varphi$ 是函数 $A\to\R$。设 $P$ 是 $\R^n$ 中由 $A$ 的凸包给出的多面体；于是要求 $\varphi$ 是 $P$ 上某个凸分段线性函数的限制，而后者的线性性极大区域给出 $P$ 的一个多面体分解，其顶点恰为 $A$ 中的点。

\transnote{易错预警：$T^{\varphi(\alpha)}$ 中的 $T$ 是一个小实数，充当退化参数（相当于 SYZ 图景中的尺度）；$\varphi$ 才是分段线性函数，两者不要与温度、周期或平移混淆。取 $T\to0$ 时，各项按 $\varphi(\alpha)$ 的大小分层，这正是热带极限的来源。}

% ===== 原书 p.16 =====
我们将进一步要求 $f$ 满足以下条件：
\begin{enumerate}[label=\arabic*.,leftmargin=2em,itemsep=0.2em]
\item $A$ 的所有非零点都落在 $P$ 的边界上；
\item 由 $\varphi$ 诱导的 $P$ 的多面体分解是极大的（分解的所有胞腔在 $\GL_n(\Z)$ 作用下都与标准单纯形合同）；
\item Laurent 多项式 $f$ 有常数项，即 $0\in A$，且 $0$ 是分解中每个极大胞腔的顶点。
\end{enumerate}

\transnote{含糊点拨：条件 (2) 中的“极大”（maximal）不是“胞腔数目最多”，而是指该多面体分解在加细意义下不能再细分：每个胞腔都是与标准单纯形合同的整多面体。条件 (1) 则把 $f$ 的支撑限制在 $P$ 的边界上，这保证 $H$ 在 $(\C^*)^n$ 中的紧化只沿边界方向退化。}

上文中，多面体 $P$ 的边界带有一个多面体分解；令 $\Sigma_X$ 为由该分解各面所张的锥构成的环簇扇。此扇的射线由 $A\subseteq\Z^n$ 中的向量生成。设 $X$ 是相应的复环簇；令 $i_D\colon D\to X$ 为环簇反典范除子的包含，该除子是反典范丛 $\cK_X^{-1}$ 的一个截面 $s$ 的零点轨迹。我们用 $\cK_X$ 记 $X$ 上作为层的典范丛，并用 $K_X$ 记典范丛的全空间；$K_X$ 也是环簇，其扇 $\tilde\Sigma_X\subseteq\R^{n+1}$ 的射线由 $(1,-\alpha)$（$\alpha\in A$）生成，且 $(1,-\alpha_i)$ 张成一个锥，当且仅当 $\alpha_i$ 张成由 $\varphi$ 诱导的多面体 $P$ 的分解中的一个胞腔。此外，令 $\tilde S$ 记 $X\times\mathbb A^1$ 内的环簇反典范除子，并令 $\hat Z$ 记 $K_X\times\mathbb A^1$ 内的环簇反典范除子（即典范丛的典范丛中的“大反典范除子”）。再令 $X\times\mathbb A^1[-1]$ 记导出概型，其中 $\mathbb A^1[-1]=\operatorname{Spec}\C[t]$，$t$ 的度为 $-1$。我们将用 $\tilde Z$ 记把 $X\times\{0\}$ 沿 $D$ 粘到 $D\times\mathbb A^1[-1]$ 上所得的导出概型。

\transnote{易错预警：$\cK_X$ 与 $K_X$ 只差一个花体，含义却不同——花体记 $X$ 上\emph{作为层}的典范丛，直立体记典范丛的\emph{全空间}（一个环簇）。原文正是用这一对记号区分层与全空间，译文中不可合并。}

\transnote{易错预警：$\hat Z$ 与 $\tilde Z$ 形状相近但不同。$\hat Z$ 是 $K_X\times\mathbb A^1$ 内的环簇反典范除子；$\tilde Z$ 是把 $X\times\{0\}$ 沿 $D$ 粘到 $D\times\mathbb A^1[-1]$ 上得到的导出概型。前者是除子，后者是导出概型，粘合的方向也容易读反。}

% ===== 原书 p.17 =====
在上述条件 (3) 下，$X$ 必是光滑的，因此 $\cK_X^{-1}$ 的所有非零截面都是正则的。

假设 (1) 与 (2) 将保证镜像的环簇 $X$ 是光滑的。此外，热带阿米巴 $\Pi_f\subseteq\R^n$ 的补集将恰有一个紧分支（由条件 (3)），于是 $z_0f+1$ 的镜像可典范地等同于典范丛 $K_X$；并且取 $T$ 充分小，可使 $f$ 的全部临界点（由条件 (2)）落在 $1$ 的任意小邻域上方 [Aur18]。若愿意在 [GS22] 的环簇叠框架下工作，假设 (1) 与 (2) 可以减弱：我们将在下文对此作说明。就下文的目的而言，我们将取 $\varphi$ 为函数 $\varphi(\alpha)=1$（对每个 $\alpha\in A\setminus\{0\}$）的一个一般小扰动。

\begin{remarkx}\label{rem:stacky}
条件 (2) 可以减弱为只要求多面体分解的胞腔是单纯形，而不必具有最小体积。给定扇 $\Sigma_X$ 以及每条射线的生成元 $\alpha$ 的选择，我们可以联系一个环簇（Deligne–Mumford）叠 $X$ [BCS05]。单纯形假设保证该叠 $X$ 是光滑的。[GS22] 的同调镜像对称定理正是在这一条件下陈述的，本论文的全部结果在叠的情形下仍成立。在此较弱的条件下，$X$ 作为 Deligne–Mumford 叠是光滑的 [GS22]。本论文所用的一切代数几何结果，在此情形下应当同样适用。
\end{remarkx}

$H$ 的补集可用两种自然的方式看作 $(\C^*)^{n+1}$ 中的极仿射超曲面。它可以由 Laurent 多项式 $\tilde f_0=f+z_0^{-1}$ 定义，也可以由 $\tilde f_\infty=z_0f+1$ 定义，其中 $z_0\in\C^*$ 是 $(\C^*)^{n+1}$ 的第一个因子。当然，作为代数簇二者是相同的，我们将它们统称为 $(\C^*)^n\setminus H$。

\transnote{含糊点拨：$\tilde f_0$ 与 $\tilde f_\infty$ 是同一个补集在 $(\C^*)^{n+1}$ 中的两种极仿射描述，区别只在新增的 $z_0$ 方向上的行为：前者在 $z_0=0$ 处有极点，后者在 $z_0=\infty$ 处有极点。二者作为代数簇相同，但下面的分次数据不同，这正是 §1.5 后半部分的要点。}

% ===== 原书 p.18 =====
考虑 $\phi\colon(\C^*)^n\to\R$，$\phi(z)=|\operatorname{Log}(z)|^2$，即 $(\C^*)^n$ 上的标准 Stein 函数；通过限制，$H$ 与 $(\C^*)^n\setminus H\subseteq(\C^*)^{n+1}$ 都具有 Stein 流形的结构，从而是 Liouville 流形，这使我们能够谈论它们的 wrapped Fukaya 范畴。注意，把 $\phi$ 限制到 $(\C^*)^n\setminus H$ 所得的 Stein 结构，与 [Jef22, p.13] 中在 Stein 流形里超曲面的补集上构造的 Stein 结构是简单 Liouville 同伦的。

对我们而言，空间 $\{\tilde f_0=0\}$ 与 $\{\tilde f_\infty=0\}$ 之间的差别，将在于我们赋予它们的分次/Maslov 数据不同——即典范线丛的一个平凡化选择，它确定深谷范畴上的一个 $\Z$-分次。给定超曲面 $H\subseteq(\C^*)^n$，满足 $H=\{f=0\}$ 的函数 $f$ 给出 $H$ 的典范丛的一个平凡化 $\eta$：把
\[
\Omega=\frac{\dd z_1}{z_1}\wedge\cdots\wedge\frac{\dd z_n}{z_n}
\]
写作 $H$ 上的标准全纯最高阶形式，我们便可通过条件 $\eta\wedge\dd f=\Omega$ 来定义 $\eta$。类似地，$\dd\tilde f_0$ 与 $\dd\tilde f_\infty$ 给出 $(\C^*)^n\setminus H$ 的典范丛的平凡化 $\eta_0,\eta_\infty$。具体地，$\eta_0$ 由限制 $\Omega|_{(\C^*)^n\setminus H}$ 得到，而 $\eta_\infty$ 将是 $f^{-1}\Omega|_{(\C^*)^n\setminus H}$。注意 $\eta_0$ 因此可以延拓越过被去掉的纤维 $H$，而 $\eta_\infty$ 不能。把分次从 $\eta_0$ 换到 $\eta_\infty$，相当于用 $f$ 去除 $\eta_0$。就上同调而言，这对应于用类 $[\arg(f)]\in H^1((\C^*)^n\setminus H;\Z)$ 来扭平凡化 $\eta_0$。

这两个平凡化定义了两个\emph{不同的} dg 范畴，我们分别记之为 $\cW_0((\C^*)^n\setminus H)$ 与 $\cW_\infty((\C^*)^n\setminus H)$。注意它们有\emph{不同的对象}：$(\C^*)^n\setminus H$ 中并非每个拉格朗日量对两种分次都有分次提升。或者，我们可以考虑一个\emph{分次 $A_\infty$ 范畴} $\cW_{0,\infty}((\C^*)^n\setminus H)$，其对象进一步限制为那些对两种分次都容许分次提升者。

\transnote{背景直觉：$\eta_0$ 与 $\eta_\infty$ 只差乘以 $f$ 的幂，但前者在 $H$ 上无极点、后者有极点，于是“哪些拉格朗日量能带分次提升”就不同。这正是同一个辛流形上出现两个不同 dg 范畴的机制；$[\arg(f)]$ 这个上同调类则度量了两种分次之间的扭结。}

% ===== 原书 p.19 =====
当然，还可以考虑 $(\C^*)^n\setminus H$ 上的其他分次结构，但这些结构从 B 侧的观点看不太自然。

\begin{remarkx}\label{rem:framing}
需要注意的是：选取 Laurent 多项式 $f$，就为 $H$ 提供了 $(\C^*)^n$ 内部的一个标架（framing）。
\end{remarkx}

我们将使用如下术语：Liouville 流形 $X$ 的\emph{扇形分解}（sectorial decomposition）或\emph{扇形粘合}（sectorial gluing），是指 $X$ 中的一个扇形超曲面 $F$，使得 $X$ 是两个 Liouville 扇形 $X_1,X_2\subseteq X$ 的并，且二者沿其边界相交，$\partial X_1=\partial X_2=F$。

下文中，所有 dg/$A_\infty$ 范畴都取为 $\C$-线性的（照常带有其内部/上同调 $\Z$-分次）；所谓\emph{分次 dg 范畴}（或 $A_\infty$ 范畴），指带有\emph{额外} $\Z$-分次的 dg 范畴；换言之，dg 范畴上的一个分次，就是一个 $\operatorname{Coh}(B\mathbb G_m)$-线性结构。dg 范畴上的一个\emph{2-周期性结构}（2-periodicity structure）是 $\C((\beta))$-线性结构的一个选择，其中 $|\beta|=2$；\emph{2-周期 dg 范畴}则指容许 2-周期性结构的 dg 范畴。等价地，把该 $\C((\beta))$-线性 dg 范畴在自然变换 $\beta$ 处局部化，便得到一个微分 $\Z/2$-分次范畴。对 $A_\infty$ 范畴，把同调分次模 $2$ 约化，就得到 $\C((\beta))$-线性的 $A_\infty$ 范畴；我们将用横杠表示同调分次模 $2$ 的约化：例如 $\overline{\cW}_0((\C^*)^n\setminus H)$ 与 $\overline{\operatorname{Coh}}(Z)$。$A_\infty$ 范畴的等价，特别会通过把分次模 $2$ 约化而下降到 $\C((\beta))$-线性 $A_\infty$ 范畴的层面。

\transnote{易错预警：$\overline{\cW}_0$、$\overline{\operatorname{Coh}}(Z)$ 上的横杠表示“同调分次模 $2$ 约化”，不是复共轭，也不是闭包或 Zariski 闭包。约化之后分次只到 $\Z/2$，这些范畴才成为 $\C((\beta))$-线性的。}

下文中，对簇 $X$，$\operatorname{Coh}(X,0)$ 将记零的矩阵分解，即相干层范畴 $\operatorname{Coh}(X)$ 的一个 $\C((\beta))$-线性版本。正如我们将在 §6.1 中解释的，可以把 $\C$-线性 dg 范畴 $\operatorname{Coh}(X)$ 作为分次零矩阵分解的范畴恢复出来；我们将后一范畴记作 $\operatorname{Coh}_{\mathbb G_m}(X,0)$。这个 $\mathbb G_m$ 作用的几何意义将在 §6.1 中解释。

% ===== 原书 p.20 =====
读者可参阅附录 A 中关于概型间态射的记号。此后，除非另有说明，下文所有函子都将是 $\Z$-分次范畴之间的 $\Z$-分次函子（即 dg 范畴之间的 dg 函子）。全文中 $\cW$ 将记 [GPS20a] 意义下（部分）wrapped Fukaya 范畴上扭复形的 $A_\infty$ 范畴；$\operatorname{Coh}$ 将记相干层的 dg 导出范畴。所有函子都取导出函子，所有余极限都应在适当的同伦意义下理解。

\subsubsection{函子式镜像对称}

对我们来说最重要的函子式镜像对称示例，是环簇中反典范除子的镜像对称：据多位作者的工作 [Abo06; Abo09; FLTZ12; Han19; Kuw20; Zho19]，对环簇 $X$（满足 §1.5 上文所列的某些假设），存在一个 Landau–Ginzburg 镜像 $((\C^*)^n,f)$，使得有拟等价
\begin{equation}\label{eq:lvl-hms}
\cW((\C^*)^n,f)\simeq\operatorname{Coh}(X),
\end{equation}
这里 $\cW$ 表示部分缠绕深谷范畴，$f$ 是 Laurent 多项式。此外，[GS22; Zho20] 用微局部层方法证明了：极仿射超曲面 $H=f^{-1}(0)$ 镜像于环簇反典范除子 $i_D\colon D\hookrightarrow X$。结合 [GPS20c] 中微局部层与 Fukaya 范畴的比较定理，这给出等价
\[
\cW(H)\simeq\operatorname{Coh}(D).
\]
此外，有交换图表 [GPS20c, Example 7.25]：
\begin{equation}\label{eq:cup-square}
\begin{tikzcd}
\cW(H) \arrow[r,"\sim"] \arrow[d,"\cup"] & \operatorname{Coh}(D) \arrow[d,"i_{D,*}"] \\
\cW((\C^*)^n,f) \arrow[r,"\sim"] & \operatorname{Coh}(X),
\end{tikzcd}
\end{equation}
其中 $\cup$ 表示 Abouzaid–Ganatra 的 cup 函子 [GPS20c]，底部的等价是来自 [Kuw20, Theorem 1.2] 并结合作者 [GPS20c] 的工作所得的等价 (1.1)。这是我们证明的各定理的原型：我们证明同调镜像对称尊重镜像对称两侧某些自然定义的函子。

\transnote{背景直觉：交换图表 (1.2) 是全文的“原型”。右侧 $i_{D,*}$ 是除子 $D\hookrightarrow X$ 的推前；左侧的 $\cup$ 是它在 A 侧的对应物——把定义在 $H$（$D$ 的镜像）上的对象“延拓”到 $(\C^*)^n\setminus H$ 上。所谓函子式镜像对称，就是断言这类由几何自然给出的函子，在镜像等价下逐一对得上。}

% ===== 原书 p.21 =====
Auroux 在 [Aur18] 中观察到，补集 $(\C^*)^n\setminus H$ 也是一个极仿射超曲面，镜像于典范丛 $K_X$ 内的环簇反典范除子 $Z$。我们证明的函子性结果将表明，[GS22] 的同调镜像对称等价把往来于 $\cW((\C^*)^n\setminus H)$ 的、由辛几何定义的函子，与来自镜像 $Z$ 的代数几何的各种拉回、推前映射编织在一起。我们将非正式地说：当存在形如 (1.2) 的交换图表时（相对于指定的镜像对称等价），这些函子在\emph{镜像对称下相对应}。这些结果的精确陈述在下文 §1.6 中给出。
